\DocumentMetadata{
  pdfversion=2.0,
  pdfstandard=ua-2,
  lang=en-US,
  tagging=on,
  tagging-setup={math/setup=mathml-SE,tabsorder=structure}
}
\documentclass[11pt,letterpaper]{article}
\usepackage[margin=0.68in,includefoot]{geometry}
\usepackage{newcomputermodern}
\usepackage{amsmath,amscd,mathtools}
\usepackage{tikz,adjustbox}
\providecommand{\square}{\mdlgwhtsquare}
\usepackage[hidelinks]{hyperref}
\hypersetup{
  pdftitle={Analysis First-Year Exam, January 2024, Part 1},
  pdfauthor={University of Florida Department of Mathematics},
  pdfsubject={Graduate mathematics examination},
  pdfdisplaydoctitle=true,
  bookmarksopen=true
}
\setcounter{secnumdepth}{0}
\setlength{\parindent}{0pt}
\setlength{\parskip}{0.28em}
\setlength{\emergencystretch}{3em}
\setlength{\leftmargini}{2.15em}
\setlength{\leftmarginii}{2.65em}
\setlength{\leftmarginiii}{3em}
\pagestyle{empty}
\raggedbottom
\allowdisplaybreaks
\NewTaggingSocketPlug{block/list/label}{examproblem}{%
  \tagstructbegin{tag=Lbl}\tagstructend%
  \tagstructbegin{tag=\LBody}%
  \tagstructbegin{tag=\ProblemHeadingTag,title-o={\ProblemHeadingText}}%
  \tagmcbegin{tag=\ProblemHeadingTag,actualtext={\ProblemHeadingText}}%
  #2%
  \tagmcend%
  \tagstructend%
}
\newenvironment{examproblems}[2]{%
  \def\ProblemHeadingTag{#2}%
  \AssignTaggingSocketPlug{block/list/label}{examproblem}%
  \begin{enumerate}%
  \setcounter{enumi}{#1}%
  \renewcommand{\labelenumi}{\textbf{\theenumi.}}%
  \setlength{\topsep}{0.35em}%
  \setlength{\partopsep}{0pt}%
  \setlength{\itemsep}{0.48em}%
  \setlength{\parsep}{0.18em}%
}{\end{enumerate}}
\newenvironment{examsubparts}{%
  \AssignTaggingSocketPlug{block/list/label}{default}%
  \begin{enumerate}%
  \setlength{\topsep}{0.18em}%
  \setlength{\partopsep}{0pt}%
  \setlength{\itemsep}{0.16em}%
  \setlength{\parsep}{0.1em}%
}{\end{enumerate}}
\newenvironment{examinstructions}{%
  \par\begingroup\itshape
}{%
  \par\endgroup\vspace{0.18em}
}
\makeatletter
\renewcommand\subsection{\@startsection{subsection}{2}{0pt}%
  {-1.25ex plus -.25ex minus -.1ex}%
  {0.55ex plus .1ex}%
  {\normalfont\large\bfseries}}
\renewcommand{\@maketitle}{%
  \newpage\null\vskip -1em%
  {\footnotesize\bfseries
    \href{https://www.ufl.edu/}{University of Florida}, \href{https://math.ufl.edu/}{Department of Mathematics}\par}%
  \vskip -0.5em%
  \tagstructbegin{tag=H1,title={Analysis First-Year Exam, January 2024, Part 1}}%
  \tagpdfparaOff%
  \tagmcbegin{tag=H1}%
    {\LARGE\bfseries\href{https://gma.math.ufl.edu/exams/analysis-first-year/}{Analysis First-Year Exam}, January 2024, Part 1\par}%
  \tagmcend%
  \tagpdfparaOn%
  \tagstructend%
  \vskip 0.25em%
  \tagmcbegin{artifact}%
    \hrule height 1pt%
  \tagmcend%
  \vskip 1em%
}
\makeatother
\title{Analysis First-Year Exam, January 2024, Part 1}
\author{}
\date{}
\begin{document}
\maketitle
\thispagestyle{empty}
\begin{examinstructions}
Answer FOUR questions in detail. State carefully any results used without proof.
\end{examinstructions}
\begin{examproblems}{0}{H2}
\def\ProblemHeadingText{Problem 1.}
\item
Let \(f:(0,\infty)\to\mathbb{R}\) be continuous. Prove that if
\[
\limsup_{n\to\infty} f(n)>0>\liminf_{n\to\infty} f(n)
\]
 then~\(f\) has infinitely many zeros. Does the same conclusion hold if one of the strict inequalities is replaced by a weak inequality?
\def\ProblemHeadingText{Problem 2.}
\item
Let \(f:\mathbb{R}\to\mathbb{R}\) be continuous and assume that \(f^{-1}(y)\subseteq\mathbb{Q}\) whenever \(y\in\mathbb{R}\setminus\mathbb{Q}\). What can be deduced about~\(f\)?
\def\ProblemHeadingText{Problem 3.}
\item
Let the function \(f:(0,1)\to\mathbb{R}\) be uniformly continuous.
\begin{examsubparts}
\item[\textnormal{(i)}]
Prove that~\(f\) extends to a continuous function \(F:[0,1]\to\mathbb{R}\).
\item[\textnormal{(ii)}]
Does the same conclusion follow if~\(f\) is merely continuous? What if~\(f\) is continuous and bounded? Justify your assertions.
\end{examsubparts}
\def\ProblemHeadingText{Problem 4.}
\item
Let \(f:X\to Y\) be a map between metric spaces and consider the following two statements. For each implication (i)~\(\Rightarrow\) (ii) and (ii)~\(\Rightarrow\) (i), give a proof (if true) or a counterexample (if false).
\begin{examsubparts}
\item[\textnormal{(i)}]
if \((x_n)\) is a Cauchy sequence in~\(X\), then \((f(x_n))\) is Cauchy in~\(Y\);
\item[\textnormal{(ii)}]
\(f\) is uniformly continuous.
\end{examsubparts}
\def\ProblemHeadingText{Problem 5.}
\item
The function \(f:\mathbb{R}\to\mathbb{R}\) is differentiable, except perhaps at \(a\in\mathbb{R}\). Prove that if the real limit \(\lim_{t\to a}f'(t)\) exists then~\(f\) is in fact differentiable at~\(a\).
\end{examproblems}
\end{document}
